% GENERATED FILE. Edit canonical lessons, then run npm run generate:books.
% canonical-source-hash: fnv1a64:23b5944c5003886b
% canonical-lessons: TA-C67-shirt, TA-C67-veshti, TA-W23-read-sattai, TA-C67-sari, TA-C67-sandal, TA-C67-cap

\chapter{What You Wear}
\label{ch:ta-what-you-wear}
\hlchaptermodality{67}

\begin{chapteropening}
\textbf{By the end of this chapter:} \emph{I can name a shirt, a veshti, a sari, sandals and a cap in Tamil, and run the five from the head down to the floor.}

Complete TA-C67-cap and say all five words of the chapter in order.
\end{chapteropening}

\section[saṭṭai]{\ta{சட்டை} (saṭṭai) — a shirt}

\label{lesson:TA-C67-shirt}



\begin{quote}
Before the new one: one word back from last chapter: what did \emph{iḍadu} mean?
\end{quote}

\subsection*{You'll want to know: \ta{சட்டை}}
\textbf{\ta{சட்டை}} (\emph{saṭṭai}) — \textquotedblleft{}a shirt\textquotedblright{}.

The ordinary word for a shirt — the thing that goes over the top half of you, in a shop or on a washing line.

You already own \ta{துணி} (\emph{tuṇi}), cloth. The two are one step apart: \ta{துணி} is the material, and \ta{சட்டை} is the thing cut and sewn out of it. Ask for \ta{துணி} and you get a length of fabric; ask for a \ta{சட்டை} and you get a garment.

Say the middle hard. The doubled \textbf{\ta{ட்ட}} in \ta{சட்டை} is the same doubled \ta{ட} you already hear in \ta{பெட்டி} (\emph{peṭṭi}), a box — and it is what keeps this word \emph{saṭṭai} rather than a softer \emph{saḍai}.

\begin{grammarlens}[title={What you've built}]
The first of five things you can wear.
\end{grammarlens}

\subsection*{Your turn}
\begin{itemize}
\raggedright
  \item \emph{Say it:} \emph{saṭṭai}
  \item \emph{Say it:} \emph{tuṇi}, then \emph{saṭṭai} — and say which one is still just cloth
  \item \emph{Recall:} say \emph{valadu}, then say \emph{iḍadu}, then say \emph{saṭṭai}
\end{itemize}

\begin{tcolorbox}[breakable,colback=teal!4,colframe=teal!35!black,title={Before you move on}]
Say it once more, holding the doubled \ta{ட}.
\end{tcolorbox}
\section[vēṭṭi]{\ta{வேட்டி} (vēṭṭi) — a veshti, the cloth a man wraps at the waist}

\label{lesson:TA-C67-veshti}



\begin{quote}
Before the new one: what did \emph{saṭṭai} mean?
\end{quote}

\subsection*{You'll want to know: \ta{வேட்டி}}
\textbf{\ta{வேட்டி}} (\emph{vēṭṭi}) — \textquotedblleft{}a veshti, the cloth a man wraps at the waist\textquotedblright{}.

Not a sewn garment at all. A \ta{வேட்டி} is a single uncut length of \ta{துணி}, wrapped round the waist and tucked — which is why the cloth word sits so close behind it.

English borrowed the northern word \emph{dhoti} for the same idea. In Tamil Nadu the word you will hear, and the word to say, is \ta{வேட்டி}.

Its middle is the doubled \ta{ட} you met one lesson ago in \ta{சட்டை}, riding a long \ta{ே} this time.

\begin{grammarlens}[title={What you've built}]
Two garments, and the same doubled consonant in both.
\end{grammarlens}

\subsection*{Your turn}
\begin{itemize}
\raggedright
  \item \emph{Say it:} \emph{vēṭṭi}
  \item \emph{Say it:} \emph{saṭṭai}, then \emph{vēṭṭi} — and say which one was never cut
  \item \emph{Recall:} say \emph{iḍadu}, then say \emph{saṭṭai}, then say \emph{vēṭṭi}
\end{itemize}

\begin{tcolorbox}[breakable,colback=teal!4,colframe=teal!35!black,title={Before you move on}]
Which of the two is one uncut length of cloth? (\textbf{\ta{வேட்டி}}.)
\end{tcolorbox}
\section[saṭṭai]{\ta{சட்டை} — what a doubled letter looks like}

\label{lesson:TA-W23-read-sattai}



\begin{quote}
Before the page: you have said \emph{saṭṭai} and \emph{vēṭṭi} aloud. Now look at the first one on the page.
\end{quote}

\subsection*{Script you'll notice: \ta{சட்டை} — what a doubled letter looks like}
\textbf{\ta{சட்டை}} brings no new shape. It brings a familiar one used twice in a row, which is worth watching once slowly.

\noindent
\begin{tabularx}{\linewidth}{@{}>{\raggedright\arraybackslash}X>{\raggedright\arraybackslash}X>{\raggedright\arraybackslash}X@{}}
\toprule
\textbf{piece} & \textbf{} & \textbf{} \\
\midrule
\textbf{\ta{ச}} & \emph{sa} & the letter on its own, with its own built-in \emph{a} \\
\textbf{\ta{ட்}} & the \ta{ட} with a \textbf{\ta{புள்ளி}} on top & the dot that takes the vowel away, so nothing is said after it \\
\textbf{\ta{டை}} & \emph{ṭai} & the same \ta{ட} again, this time with the \textbf{\ta{ை}} sign \\
\bottomrule
\end{tabularx}

\begin{quote}
\textbf{\ta{ச} · \ta{ட்} · \ta{டை} $\to$ \ta{சட்டை}}
\end{quote}

The dot in the middle is doing the whole job. Without it you would have \ta{ச}-\ta{ட}-\ta{டை}, three open beats. With it, the first \ta{ட} closes and leans on the second, and the word comes out hard in the middle: \emph{saṭṭai}.

\subsection*{Your turn}
\begin{itemize}
\raggedright
  \item \emph{Read it:} \textbf{\ta{ச}} — then \textbf{\ta{ட்}} — then \textbf{\ta{டை}}
  \item \emph{Point to:} the \ta{புள்ளி}, and say what it takes away
  \item \emph{Read it:} \textbf{\ta{சட்டை}}, then say what it means
  \item \emph{Recall:} say \emph{saṭṭai}, then say \emph{vēṭṭi}, then read \textbf{\ta{சட்டை}}
\end{itemize}

\begin{tcolorbox}[breakable,colback=teal!4,colframe=teal!35!black,title={Before you move on}]
Read \textbf{\ta{சட்டை}}. Which mark makes the middle hard? (The \textbf{\ta{புள்ளி}} on the first \ta{ட}.)
\end{tcolorbox}
\section[puḍavai]{\ta{புடவை} (puḍavai) — a sari}

\label{lesson:TA-C67-sari}



\begin{quote}
Before the new one: read \textbf{\ta{சட்டை}} once, then say what it means.
\end{quote}

\subsection*{You'll want to know: \ta{புடவை}}
\textbf{\ta{புடவை}} (\emph{puḍavai}) — \textquotedblleft{}a sari\textquotedblright{}.

The everyday Tamil word for the draped cloth a woman wears.

English took \textbf{sari} in from the north. Tamil's own word is \ta{புடவை}, and it is the one that works in a Tamil shop: say \ta{புடவை} and you will be shown exactly what you asked for.

It ends in the \textbf{\ta{ை}} you have now met three times over — in \ta{தலை} (\emph{talai}), the head, and twice inside \ta{சட்டை}.

\begin{grammarlens}[title={What you've built}]
Three garments, and a Tamil word English never borrowed.
\end{grammarlens}

\subsection*{Your turn}
\begin{itemize}
\raggedright
  \item \emph{Say it:} \emph{puḍavai}
  \item \emph{Say it:} \emph{talai}, \emph{saṭṭai}, \emph{puḍavai} — three words, one ending
  \item \emph{Recall:} say \emph{vēṭṭi}, then read \textbf{\ta{சட்டை}}, then say \emph{puḍavai}
\end{itemize}

\begin{tcolorbox}[breakable,colback=teal!4,colframe=teal!35!black,title={Before you move on}]
Which ending do \ta{தலை}, \ta{சட்டை} and \ta{புடவை} share? (The \textbf{\ta{ை}}.)
\end{tcolorbox}
\section[seruppu]{\ta{செருப்பு} (seruppu) — a sandal, footwear}

\label{lesson:TA-C67-sandal}



\begin{quote}
Before the new one: what did \emph{puḍavai} mean?
\end{quote}

\subsection*{You'll want to know: \ta{செருப்பு}}
\textbf{\ta{செருப்பு}} (\emph{seruppu}) — \textquotedblleft{}a sandal, footwear\textquotedblright{}.

Footwear — and in Tamil it is really a doorway word.

\ta{செருப்பு} comes off at the \ta{வாசல்} (\emph{vāsal}) before anyone steps \ta{உள்ளே}. So the thing you have just named is the one thing that stays \ta{வெளியே}: three words from earlier chapters, and this one is the reason they get used every day.

\begin{grammarlens}[title={What you've built}]
Four, and the run has reached the floor.
\end{grammarlens}

\subsection*{Your turn}
\begin{itemize}
\raggedright
  \item \emph{Say it:} \emph{seruppu}
  \item \emph{Say it:} \emph{seruppu}, then \emph{vāsal} — and say which side of the door each one is on
  \item \emph{Recall:} read \textbf{\ta{சட்டை}}, then say \emph{puḍavai}, then say \emph{seruppu}
\end{itemize}

\begin{tcolorbox}[breakable,colback=teal!4,colframe=teal!35!black,title={Before you move on}]
Where does the \ta{செருப்பு} stay — \ta{உள்ளே} or \ta{வெளியே}? (\textbf{\ta{வெளியே}}.)
\end{tcolorbox}
\section[toppi]{\ta{தொப்பி} (toppi) — a cap, a hat}

\label{lesson:TA-C67-cap}



\begin{quote}
Before the new one: what did \emph{seruppu} mean, and what did \emph{puḍavai} mean?
\end{quote}

\subsection*{You'll want to know: \ta{தொப்பி}}
\textbf{\ta{தொப்பி}} (\emph{toppi}) — \textquotedblleft{}a cap, a hat\textquotedblright{}.

The last of the five, and the highest one: a \ta{தொப்பி} sits \ta{மேலே}, on the \ta{தலை}.

Its \textbf{\ta{ொ}} is the short o sign you met long before this chapter, and its doubled \textbf{\ta{ப்ப}} works exactly like the doubled \ta{ட} you took apart in \ta{சட்டை} — a \ta{புள்ளி} closing the first letter onto the second.

With this one the run has both ends: \ta{தொப்பி} at the top, \ta{செருப்பு} at the floor.

\begin{grammarlens}[title={What you've built}]
Five things to wear, top to bottom.
\end{grammarlens}

\subsection*{Your turn}
\begin{itemize}
\raggedright
  \item \emph{Say it:} \emph{toppi}
  \item \emph{Say it:} all five from the head down — \emph{toppi}, \emph{saṭṭai}, \emph{vēṭṭi}, \emph{puḍavai}, \emph{seruppu}
  \item \emph{Recall:} say \emph{puḍavai}, then say \emph{seruppu}, then say \emph{toppi}
\end{itemize}

\begin{tcolorbox}[breakable,colback=teal!4,colframe=teal!35!black,title={Before you move on}]
Which of the five sits highest, and which stays \ta{வெளியே}? (\textbf{\ta{தொப்பி}}; \textbf{\ta{செருப்பு}}.)
\end{tcolorbox}
